% Version: v3 (2026-11-01), corrections from a review of v2.
% v3 changes, all local: the abstract now carries the 2-torsion-free hypothesis in the two places
% where the corresponding theorem has it (the commutative case and the vanishing on eigenvectors);
% the standing hypothesis of Section 2 is stated as enrichment over abelian groups instead of "every
% set of morphisms is an abelian group"; Definition 2.1 is glossed (K, S, Sweedler notation) and the
% sentence relating mu_- to the braided commutator now assumes a braided algebra, as in Ion's
% Definition 1.4; Definition 4.5 writes the second morphism as zeta(l-hat o ev) with l-hat = zeta(l),
% as in Barnes-Schenkel-Szabo Definition 3.4 (l o ev is not composable); the kernel diagram of
% Theorem 4.3 was wrong (it drew a zero map as a monomorphism and omitted Der(mu_+)) and is redrawn,
% and the display of Theorem 4.3 is split to remove an overfull box; Remark 4.4 cites the sentence
% after Pareigis' Definition 5.5, where the object Der(A) is actually introduced; the introduction no
% longer claims that the proofs use only naturality and addition; and Section 2 records the
% specialisation to symmetric monoidal categories (mu(1 +- c) = mu +- mu c, (1 + c)(1 - c) = 0, and
% the twist of an involution is an involution).
% v3, second review pass (reference-grounded): the abstract states the standing hypothesis as
% enrichment over abelian groups, as Section 2 does; Definition 4.5 restores the hypotheses of
% Barnes-Schenkel-Szabo Definition 3.4 (a quasitriangular quasi-Hopf algebra and a braided
% commutative algebra in its module category); the sentence introducing Ion's braided commutator
% names his Definition 1.4; the Fuchs-Runkel-Schweigert map is renamed (their sigma is an
% endomorphism of A, not of A(x)A); Claim 2.14 of Graziani et al. is reported as stated there, not
% proved; and the paragraph after the proof of Theorem 4.3 records that only Lemma 4.2 uses the
% symmetry, so that Der(mu) is available in a braided closed category with the kernel of Phi.
% The kernel diagram after Theorem 4.3 now labels its middle vertical by id: the label Phi_+ that the
% previous pass put there is a morphism [A,A] -> [A(x)A,A], not an endomorphism of [A,A].
% Section 2 no longer claims that every statement holds in any symmetric monoidal category without the
% further hypotheses (2-torsion freeness, closedness, kernels), nor that the section contains no new
% results (it carries the symmetric specialisation); the abstract says "every case treated here"
% instead of "all cases"; and "symmetric algebra" in the sense of Graziani et al. is glossed.
% In Section 4, the first morphism of the Barnes-Schenkel-Szabo equalizer is now
% zeta([.,.] o (id(x)l-hat)) with [.,.] the internal commutator of end(A), as their (3.10)-(3.11)
% have it; "closed" is defined before Theorem 4.3; Remark 4.4 drops the symmetry (it is not used) and
% gives the one-line argument from Pareigis' Corollary 5.6; and the family defect is identified with
% Pareigis' rule at the universal family T = end(A), f = ev, not at the unit object.
% The two remaining claim-ok comments in Section 5 now say Theorem 5.1, and the Pareigis entry of the
% bibliography carries the journal data (Appl. Categ. Structures 6 (1998) 151-175), the archived copy
% being the journal version.
% Theorem 5.2 now carries the hypotheses of its source (associative, characteristic not 2), semiprime
% is defined, and Cusack's Corollary 5 is named as the first proof, as the survey does; the sentence
% introducing Bresar's Theorem 1 names its hypothesis instead of saying "a suitable hypothesis".
% Version: v2 (2026-10-30), the antisymmetrized product mu(1 - c) is the braided commutator of Ion, and commutative means c-commutative; Ion 2010 cited with Definitions 1.4 and 1.5.
\documentclass[]{interact}
\usepackage[utf8]{inputenc}

\usepackage{mathtools}
\usepackage{amsmath,amssymb,amsthm}
\usepackage{amsfonts}
\usepackage{enumerate}
\usepackage{enumitem}

\usepackage[hidelinks,pagebackref=false]{hyperref}
\usepackage[all]{hypcap}
\hypersetup{linktoc=all}

\usepackage{tikz}
\usepackage{tikz-cd}
\usepackage[yyyymmdd,hhmmss]{datetime}
\usepackage{float}

\usepackage[numbers,sort&compress]{natbib}
\bibpunct[, ]{[}{]}{,}{n}{,}{,}
\renewcommand\bibfont{\fontsize{10}{12}\selectfont}
\makeatletter
\def\NAT@def@citea{\def\@citea{\NAT@separator}}
\makeatother

\theoremstyle{plain}
\newtheorem{theorem}{Theorem}[section]
\newtheorem{lemma}[theorem]{Lemma}
\newtheorem{corollary}[theorem]{Corollary}
\newtheorem{proposition}[theorem]{Proposition}
\newtheorem{claim}[theorem]{Claim}

\theoremstyle{definition}
\newtheorem{definition}[theorem]{Definition}
\newtheorem{example}[theorem]{Example}

\theoremstyle{remark}
\newtheorem{remark}[theorem]{Remark}
\newtheorem{notation}{Notation}

\allowdisplaybreaks
\newcommand{\Der}{\mathrm{Der}}
\newcommand{\Hom}{\operatorname{Hom}}
\newcommand{\id}{\mathrm{id}}

\newcommand{\paperversion}{3}

\begin{document}
\title{Derivations of the Symmetrized Product in a Braided Monoidal Category}
\author{
\name{Abderrahim Adrabi\textsuperscript{a,b}\thanks{CONTACT A. Adrabi. Email: abderrahim.adrabi@um5r.ac.ma}}
	\affil{\textsuperscript{a}Faculty of Sciences, Mohammed V University in Rabat, Morocco;\\ \textsuperscript{b}Faculty of Educational Sciences, Mohammed V University in Rabat, Morocco}
}

\history{Preprint version \paperversion, compiled \today}

\maketitle

\begin{abstract}
Let $A$ be an object of a braided monoidal category enriched over abelian groups, and let $\mu \colon A \otimes A \to A$ be a product on $A$.
Let $c$ be the braiding on $A \otimes A$.
The symmetrized product of $A$ is $\mu(1 + c)$, and the antisymmetrized product is $\mu(1 - c)$.
A derivation of a product is a morphism $D \colon A \to A$ that satisfies the Leibniz rule for that product.
The Leibniz defect of $D$ decides every case treated here.
It shows that $D$ is a derivation of $\mu(1 \pm c)$ exactly when the defect composed with $1 \pm c$ vanishes.
If no nonzero morphism is killed by $2$, a derivation of both products is a derivation of $\mu$, and for a commutative product the derivations of $\mu(1 + c)$ and of $\mu$ coincide.
In a symmetric closed category with kernels the derivations form an object that contains the derivation object of Pareigis.
Let $\alpha$ be an endomorphism of $A$ with $\alpha^2 = \id_A$ and let $\sigma = (\alpha \otimes \alpha)c$ be the twist.
If $\mu$ is commutative and $D$ commutes with $\alpha$, then a derivation of $\mu_\sigma$ has a defect that is killed by $1 + (\alpha \otimes \alpha)$; if moreover no nonzero morphism is killed by $2$, it vanishes on pairs of eigenvectors of $\alpha$ with the same eigenvalue.
\end{abstract}

\begin{keywords}
braided monoidal category; derivation; symmetrized product; Jordan derivation; Leibniz defect
\end{keywords}

\begin{amscode}
MSC 2020: 18M15; 16W25; 17C50
\end{amscode}

\section{Introduction}\label{sec:intro}

A derivation of an algebra $A$ is a linear map $D$ with $D(xy) = D(x)y + xD(y)$.
\citet[Part I]{Jacobson1937} studied the derivations of an arbitrary nonassociative algebra.
Every associative algebra carries a second product, the sum $x \circ y = xy + yx$.
Jordan, von Neumann and Wigner form a product $ab$ from the two products $a \times b$ and $b \times a$ of an associative system, and call this quasi-multiplication \citep[p.~64]{JordanVonNeumannWigner1934}.
\citet[p.~1104]{Herstein1957} calls a derivation of this product a Jordan derivation.
The sum is twice the Jordan product $\tfrac12 (xy + yx)$ of an associative algebra \citep[Example 2.1]{BresarZelmanov2025}.
\citet[Theorem 3.1]{Herstein1957} proved that every Jordan derivation of a prime ring of characteristic not $2$ is a derivation.

In a braided monoidal category the braiding $c$ replaces the swap of vector spaces \citep[Definition 2.1]{JoyalStreet1993}.
The product $\mu(1 + c)$ is the sum of the product and its braided opposite. % claim-ok: the definition of mu_+ given in Section 2
\citet[Claim 2.14]{GrazianiMakhloufMeniniPanaite2015} state that, for a symmetric algebra $(A, \mu, u, c)$ in their sense (a braided unital algebra with $c^2 = \id_A$), the morphism $\mu \circ (\mathrm{Id}_{A \otimes A} - c)$ defines a braided Lie algebra structure on $A$; the claim is recalled there from Goyvaerts and Vercruysse.
\citet[Definition 5.5]{Pareigis1998} defines derivations of $\mu$ in the category of Yetter--Drinfeld modules.
We found no source that treats derivations of $\mu(1 + c)$ for a general braided monoidal category.

This paper is the third of a series on symmetrized products. % claim-ok: states the place of this paper in the series
The first paper defines braided-symmetric derivations of an algebra, which are derivations of the product $\mu + \mu\sigma$ for an involution $\sigma$ of $A \otimes A$ \citep[Definition 2.1]{Adrabi2026a}.
The second treats the Lie derivations of an associative algebra as the derivations of its antisymmetrized product \citep[Section 4]{Adrabi2026b}.
A Lie derivation of an associative algebra is a linear map $D$ with $D([x,y]) = [D(x), y] + [x, D(y)]$ \citep[p.~183]{Martindale1964}.
Here the category is general, and the proofs use the axioms of a braided monoidal category and the additivity of morphisms, together with, where a statement says so, symmetry, $2$-torsion freeness, or closedness and kernels.

Section~\ref{sec:prelim} recalls the notions of the sources.
Section~\ref{sec:defect} proves that the Leibniz defect decides the derivations of $\mu(1 \pm c)$ (Proposition~\ref{prop:defect-sym}) and that a derivation of both products is a derivation of $\mu$ (Theorem~\ref{thm:both}).
Section~\ref{sec:commutative} treats a commutative product (Theorem~\ref{thm:commutative}) and the derivation object (Theorem~\ref{thm:object}).
Section~\ref{sec:twisted} twists the braiding by a morphism $\alpha$ (Theorem~\ref{thm:twisted}).
We end with one open question.

\section{Preliminaries}\label{sec:prelim}
This section recalls the notions of the sources that the paper uses; apart from the specialisation to symmetric monoidal categories recorded below, it contains no new results.

Throughout, $\mathcal{C}$ is a braided monoidal category enriched over abelian groups: every hom-set is an abelian group, composition is bilinear, and the tensor product of morphisms is bilinear.
Write $c_{X,Y} \colon X \otimes Y \to Y \otimes X$ for the braiding. % claim-ok: states the setting of the paper
We omit the isomorphisms of associativity and of the unit, by the coherence theorem for braided monoidal categories \citep[Theorem 2.5]{JoyalStreet1993}.
Let $\mu \colon A \otimes A \to A$ be a morphism, and write $c = c_{A,A}$.
We do not assume that $\mu$ is associative.
Every statement below holds, in particular, in a symmetric monoidal category, in which $c_{Y,X} \circ c_{X,Y} = \id$ for all $X, Y$; the hypotheses that go beyond braidedness are named in the statements that use them, namely $2$-torsion freeness in Theorems~\ref{thm:both}, \ref{thm:commutative} and \ref{thm:twisted}(iii), and closedness with kernels in Theorem~\ref{thm:object}.
There the braiding on $A \otimes A$ is an involution, so that $\mu(1 \pm c) = \mu \pm \mu c$ and
\[
(1 + c)(1 - c) = (1 - c)(1 + c) = 0,
\]
and the twist $\sigma = (\alpha \otimes \alpha)c$ of Section~\ref{sec:twisted} is an involution whenever $\alpha$ is.

We recall Pareigis' definition of a derivation, which is the model for the notion below. Here $\mathsf{C}$ is a category of Yetter--Drinfeld modules.

\begin{definition}[{\citep[Definition 5.5]{Pareigis1998}}]\label{def:pareigis}
Let $A$ be an algebra in $\mathsf{C}$ and let $\mathrm{end}(A)$ be the inner endomorphism object of $A$ in $\mathsf{C}$, that is, the Yetter--Drinfeld module $\mathrm{end}(A)$ satisfying $\mathsf{C}(X \otimes A, A) \cong \mathsf{C}(X, \mathrm{end}(A))$ for all $X \in \mathsf{C}$.
His description of it is $\mathrm{end}(A) = \{ f \in \Hom(A,A) \mid \exists \sum f_{(0)} \otimes f_{(1)} \in \Hom(A,A) \otimes K \;\forall a \in A : \sum f_{(0)}(a) \otimes f_{(1)} = \sum f(a_{(0)})_{(0)} \otimes f(a_{(0)})_{(1)} S(a_{(1)}) \}$, which is the Yetter--Drinfeld module with the required universal property.
The object $\mathrm{end}(A)$ operates on $A$ by a canonical map $\mathrm{ev} \colon \mathrm{end}(A) \otimes A \to A$ with $\mathrm{ev}(f \otimes a) = f(a)$.
A derivation from $A$ to $A$ is a linear map $d \colon A \to A$ in $\mathrm{end}(A)$ such that $d(ab) = d(a)b + (1 \otimes d)(\tau \otimes 1)(d \otimes a \otimes b)$ for all $a, b \in A$.
Here $K$ is the Hopf algebra whose Yetter--Drinfeld modules form $\mathsf{C}$, $S$ is its antipode, and $\sum f_{(0)} \otimes f_{(1)}$, $\sum a_{(0)} \otimes a_{(1)}$ are Sweedler's notation for the components of $f$ and of the coaction on $A$.
Equivalently, as \citet[proof of Corollary 5.6]{Pareigis1998} puts it, $d$ is a derivation when $m \circ (d \otimes 1 + 1 \otimes d) = d \circ m$ for the product $m$ of $A$, the tensor of $\mathrm{end}(A)$ with $A$ being the braided one, $(x \otimes y)(a \otimes b) = (\mathrm{ev} \otimes \mathrm{ev})(1 \otimes \tau \otimes 1)(x \otimes y \otimes a \otimes b)$.
Observe that in the symmetric situation this means $d(ab) = d(a)b + ad(b)$.
\end{definition}

For the unit object as parameter, the braiding $\tau$ in Definition~\ref{def:pareigis} disappears, and the rule becomes a square.

\begin{definition}\label{def:derivation} % ours: Definition 5.5 of Pareigis with the unit object as parameter and any morphism in place of the product
Let $\nu \colon A \otimes A \to A$ be a morphism.
A morphism $D \colon A \to A$ is a \emph{derivation of $\nu$} if $D \circ \nu = \nu \circ (D \otimes 1 + 1 \otimes D)$. % ours: the name for the unit object case
\end{definition}

Put $\mu_{\pm} = \mu \circ (1 \pm c)$.
We call $\mu_+$ the \emph{symmetrized product} and $\mu_-$ the \emph{antisymmetrized product} of $\mu$. % ours: the names for the braided versions
\citet[Definition 1.5]{Ion2010} defines a braided commutator on a braided algebra, in the sense of his Definition~1.4.

\begin{definition}[{\citep[Definition 1.5]{Ion2010}}]\label{def:ion}
Let $(A, \nabla, 1, c)$ be a braided algebra. The map $[\cdot,\cdot]_A \colon A^{\otimes 2} \to A$, $[\cdot,\cdot]_A = \nabla(\mathrm{id}_{A^{\otimes 2}} - c)$ is called the \emph{braided commutator} on $A$.
\end{definition}

When $(A, \mu, 1, c)$ is a braided algebra in the sense of \citet[Definition 1.4]{Ion2010}, the antisymmetrized product $\mu_-$ is the braided commutator of Definition~\ref{def:ion}.
The rule of Definition~\ref{def:derivation} for $\mu_+$ says that the following square commutes.
\[
\begin{tikzcd}[column sep=large]
A \otimes A \arrow[r, "\mu_+"] \arrow[d, "D \otimes 1 + 1 \otimes D"'] & A \arrow[d, "D"] \\
A \otimes A \arrow[r, "\mu_+"'] & A
\end{tikzcd}
\]
For vector spaces with the swap, $\mu_+$ is $x \otimes y \mapsto xy + yx$.
Its derivations are the Jordan derivations of \citet[p.~1104]{Herstein1957}, as shown in the first paper of the series \citep[Remark 2.4]{Adrabi2026a}.

The defect below measures how far $D$ is from a derivation of $\mu$. % claim-ok: signpost for the definition recalled next

\begin{definition}[{\citep[Definition 3.1]{Adrabi2026a}}]\label{def:defect-recalled}
The Leibniz defect of a linear map $D \colon A \to A$ is the bilinear map $E_D \colon A \times A \to A$ given, for $x, y \in A$, by $E_D(x,y) = D(xy) - D(x)y - xD(y)$.
\end{definition}

\section{The Leibniz defect}\label{sec:defect}
The map is classical: \citet[Section 2]{Cusack1975} writes it $d(a,b) = D(ab) - aDb - (Da)b$ and studies Jordan derivations through it.
The same formula gives the defect of a morphism $D \colon A \to A$ of $\mathcal{C}$, with $\Delta_D = D \otimes 1 + 1 \otimes D$:
\[
E_D = D \circ \mu - \mu \circ \Delta_D \colon A \otimes A \to A.
\]
So $D$ is a derivation of $\mu$ exactly when $E_D = 0$.
The braiding commutes with $\Delta_D$. % claim-ok: states the lemma proved next

\begin{lemma}\label{lem:natural}
$c \circ \Delta_D = \Delta_D \circ c$.
\end{lemma}

\begin{proof}
Since $c$ is natural, $c \circ (D \otimes 1) = (1 \otimes D) \circ c$ and $c \circ (1 \otimes D) = (D \otimes 1) \circ c$.
Adding the two equations gives the claim.
\end{proof}

The lemma is the commutative square
\[
\begin{tikzcd}[column sep=large]
A \otimes A \arrow[r, "c"] \arrow[d, "\Delta_D"'] & A \otimes A \arrow[d, "\Delta_D"] \\
A \otimes A \arrow[r, "c"'] & A \otimes A
\end{tikzcd}
\]
and it gives the main criterion.

\begin{proposition}\label{prop:defect-sym}
Let $D \colon A \to A$ be a morphism. Then
\[
D \circ \mu_{\pm} - \mu_{\pm} \circ \Delta_D = E_D \circ (1 \pm c).
\]
So $D$ is a derivation of $\mu_+$ exactly when $E_D \circ (1 + c) = 0$, and a derivation of $\mu_-$ exactly when $E_D \circ (1 - c) = 0$.
\end{proposition}

\begin{proof}
By Lemma~\ref{lem:natural}, $\Delta_D \circ (1 \pm c) = (1 \pm c) \circ \Delta_D$.
Hence
\[
D \mu (1 \pm c) - \mu (1 \pm c) \Delta_D = D \mu (1 \pm c) - \mu \Delta_D (1 \pm c) = E_D (1 \pm c). \qedhere
\]
\end{proof}

A derivation of $\mu$ has zero defect, so it satisfies both criteria.

\begin{corollary}\label{cor:inclusion}
A derivation of $\mu$ is a derivation of $\mu_+$ and of $\mu_-$.
\end{corollary}

\begin{proof}
If $E_D = 0$, then $E_D \circ (1 \pm c) = 0$, and Proposition~\ref{prop:defect-sym} applies.
\end{proof}

The converse needs one condition on $\mathcal{C}$. % claim-ok: signpost for the definition and theorem that follow
We say $\mathcal{C}$ is \emph{$2$-torsion free} if $f + f = 0$ implies $f = 0$ for every morphism $f$. % ours: the name for the condition on morphisms

Cusack proved that every Jordan derivation of a $2$-torsion free semiprime ring is a derivation \citep[Theorem 4 and Corollary 5]{Cusack1975}, and Bre\v{s}ar gave a short proof \citep[Theorem 1]{Bresar1988}.
The next theorem uses two products and no hypothesis on $A$. % claim-ok: signpost for the theorem that follows

\begin{theorem}\label{thm:both}
Let $\mathcal{C}$ be $2$-torsion free, and let $D \colon A \to A$ be a morphism that is a derivation of $\mu_+$ and of $\mu_-$. Then $E_D = 0$, that is, $D$ is a derivation of $\mu$.
\end{theorem}

\begin{proof}
By Proposition~\ref{prop:defect-sym}, $E_D (1 + c) = 0$ and $E_D (1 - c) = 0$.
Adding gives $E_D + E_D = 0$, so $E_D = 0$.
\end{proof}

\section{The commutative case and the derivation object}\label{sec:commutative}

Call $\mu$ \emph{commutative} if $\mu \circ c = \mu$, the condition that \citet[Definition 1.4]{Ion2010} calls $c$-commutative for a braided algebra.

Theorem~\ref{thm:both} needs both products.
For a commutative product the symmetrized product alone decides.

\begin{theorem}\label{thm:commutative}
Let $\mathcal{C}$ be $2$-torsion free and let $\mu$ be commutative.
Then the derivations of $\mu_+$ are the derivations of $\mu$.
\end{theorem}

\begin{proof}
By Lemma~\ref{lem:natural} and commutativity,
\[
E_D c = D \mu c - \mu \Delta_D c = D \mu - \mu c \Delta_D = D \mu - \mu \Delta_D = E_D.
\]
So $E_D (1 + c) = E_D + E_D$.
By Proposition~\ref{prop:defect-sym}, $D$ derives $\mu_+$ exactly when $E_D + E_D = 0$, that is, when $E_D = 0$.
\end{proof}

The proof is the triangle
\[
\begin{tikzcd}[column sep=large]
A \otimes A \arrow[r, "c"] \arrow[dr, "E_D"'] & A \otimes A \arrow[d, "E_D"] \\
 & A
\end{tikzcd}
\]
which commutes for a commutative product. % claim-ok: the triangle restates the first line of the proof of Theorem 4.1
For vector spaces with the swap, $\mu_+(x \otimes y) = 2xy$ for a commutative algebra.
The theorem then says that every Jordan derivation of a commutative algebra is a derivation when no nonzero element is killed by $2$. % claim-ok: restates Theorem 4.1 for vector spaces

The derivations of Pareigis form an object, not only a set. % claim-ok: signpost for the family defect and the theorem on the derivation object
We show that the same holds for $\mu_+$ when the braiding is a symmetry, that is, $c_{Y,X} \circ c_{X,Y} = \id$.
For an object $T$, a \emph{$T$-family} is a morphism $f \colon T \otimes A \to A$. % ours: the family form of a derivation, a morphism T tensor A to A as in the defining property of the object end(A) of Pareigis
Put
\[
\ell_f = f \otimes 1, \qquad
r_f = (1 \otimes f) \circ (c_{T,A} \otimes 1),
\]
as morphisms $T \otimes A \otimes A \to A \otimes A$.
Define the defect of the family by
\[
E_f = f \circ (T \otimes \mu) - \mu \circ \ell_f - \mu \circ r_f \colon T \otimes A \otimes A \to A.
\]
For the unit object $T$, the family $f$ is a morphism $A \to A$, and $E_f$ is the defect of the previous section.
For Yetter--Drinfeld modules and the universal family $T = \mathrm{end}(A)$, $f = \mathrm{ev}$, the equation $E_f = 0$ is the rule of Definition~\ref{def:pareigis}; for the unit object the braiding disappears and the rule is the untwisted Leibniz rule.
Write $E^+_f$ for the same expression with $\mu$ replaced by $\mu_+$.

\begin{lemma}\label{lem:family}
Let $\mathcal{C}$ be symmetric. For every $T$-family $f$,
\[
E^+_f = E_f \circ (T \otimes (1 + c)).
\]
\end{lemma}

\begin{proof}
We first show
\[
c \circ \ell_f = r_f \circ (T \otimes c), \qquad c \circ r_f = \ell_f \circ (T \otimes c).
\]
The first equation follows because $c$ is natural in its first variable, applied to $f$, and from the hexagon axiom for $c_{T \otimes A, A}$.
The second equation follows in the same way from the second variable and the hexagon axiom for $c_{A, T \otimes A}$.
It also uses $c_{A,T} \circ c_{T,A} = \id$, which is where the symmetry is needed.
Now
\[
\mu_+ \ell_f + \mu_+ r_f = \mu \ell_f + \mu r_f (T \otimes c) + \mu r_f + \mu \ell_f (T \otimes c) = \mu(\ell_f + r_f)\,(T \otimes (1 + c)).
\]
Also $f \circ (T \otimes \mu_+) = f \circ (T \otimes \mu) \circ (T \otimes (1 + c))$.
Subtracting gives the claim.
\end{proof}

The lemma shows that the family defect for $\mu_+$ is the family defect for $\mu$ composed with $T \otimes (1 + c)$.
This gives the objects.
Here \emph{closed} means that each $-\otimes X$ has a right adjoint $[X,-]$, so that $\Hom(T, [X,Y]) \cong \Hom(T \otimes X, Y)$.

\begin{theorem}\label{thm:object}
Let $\mathcal{C}$ be symmetric and closed, and let $\mathcal{C}$ have kernels.
Then there are objects $\Der(\mu)$ and $\Der(\mu_+)$ with
\begin{align*}
\Hom(T, \Der(\mu)) &\cong \{ f \colon T \otimes A \to A \mid E_f = 0 \},\\
\Hom(T, \Der(\mu_+)) &\cong \{ f \colon T \otimes A \to A \mid E^+_f = 0 \},
\end{align*}
natural in $T$, and $\Der(\mu)$ is a subobject of $\Der(\mu_+)$.
\end{theorem}

\begin{proof}
Write $[X,Y]$ for the internal hom, with $\Hom(T, [X,Y]) \cong \Hom(T \otimes X, Y)$.
The assignment $f \mapsto E_f$ is natural in $T$, because the braiding is natural and the tensor product is functorial in each variable.
A natural map between the functors represented by $[A,A]$ and $[A \otimes A, A]$ comes from a morphism, so $f \mapsto E_f$ is induced by a morphism $\Phi \colon [A,A] \to [A \otimes A, A]$.
By Lemma~\ref{lem:family}, the assignment $f \mapsto E^+_f$ comes from $\Phi_+ = (1 + c)^* \circ \Phi$, where $(1 + c)^*$ is composition with $1 + c$.
Put $\Der(\mu) = \ker \Phi$ and $\Der(\mu_+) = \ker \Phi_+$.
The two isomorphisms follow from the universal property of the kernel and the defining property of the internal hom.
Since $\Phi_+ = (1 + c)^* \Phi$, the kernel of $\Phi$ lies in the kernel of $\Phi_+$.
\end{proof}

The kernel diagram of the proof is
\[
\begin{tikzcd}[column sep=large]
\Der(\mu) \arrow[r, hook] \arrow[d, hook'] & {[A,A]} \arrow[r, "\Phi"] \arrow[d, "\id"'] & {[A \otimes A, A]} \arrow[d, "(1 + c)^*"] \\
\Der(\mu_+) \arrow[r, hook'] & {[A,A]} \arrow[r, "\Phi_+"'] & {[A \otimes A, A]}
\end{tikzcd}
\]
and $\Der(\mu_+) = \ker \Phi_+$ contains $\Der(\mu)$.
Only Lemma~\ref{lem:family} uses the symmetry: the morphism $\Phi$, and with it the object $\Der(\mu)$, is available already in a braided monoidal closed category in which $\Phi$ has a kernel, because the naturality of $f \mapsto E_f$ in $T$ uses only the naturality of the braiding and the bifunctoriality of $\otimes$.
The first object is the one of Pareigis. % claim-ok: restates Remark 4.4 that follows

\begin{remark}\label{rem:pareigis}
For Yetter--Drinfeld modules, with no symmetry assumed, $\Der(\mu)$ is the object $\mathrm{Der}(A)$ of \citet[Definition 5.5 and the paragraph following it]{Pareigis1998}: by the proof of his Corollary~5.6, a map $d \in \mathrm{end}(A)$ is a derivation exactly when $\mathrm m(d \otimes 1 + 1 \otimes d) = d\,\mathrm m$, an equation whose tensor is the braided one of that proof, and this is $\Phi(d) = 0$.
\end{remark}

Barnes, Schenkel and Szabo define derivations as an equalizer in the module category of a quasi-Hopf algebra \citep[Definition 3.4]{BarnesSchenkelSzabo2016}.

\begin{definition}[{\citep[Definition 3.4]{BarnesSchenkelSzabo2016}}]\label{def:bss}
Let $H$ be a quasitriangular quasi-Hopf algebra and let $A$ be a braided commutative algebra in the module category of $H$.
The \emph{derivations} of $A$ is the object $\mathrm{der}(A)$ in that module category which is defined by the equalizer
\[
\mathrm{der}(A) \longrightarrow \mathrm{end}(A) \rightrightarrows \hom(A, \mathrm{end}(A))
\]
of the two morphisms $\zeta([\cdot,\cdot] \circ (\mathrm{id} \otimes \hat l))$ and $\zeta(\hat l \circ \mathrm{ev})$, where $[\cdot,\cdot]$ is the internal commutator of $\mathrm{end}(A)$, and $\hat l = \zeta(l) \colon A \to \mathrm{end}(A)$ is the currying of the left action $l \colon A \otimes A \to A$ induced by the product of $A$.
\end{definition}

Theorem~\ref{thm:object} gives $\Der(\mu)$ as a kernel in any symmetric closed category with kernels, and the definition above gives $\mathrm{der}(A)$ as an equalizer in the module category of a quasitriangular quasi-Hopf algebra.

\section{A twisted involution}\label{sec:twisted}

Let $\alpha \colon A \to A$ be a morphism and put $\sigma = (\alpha \otimes \alpha) \circ c$ and $\mu_\sigma = \mu \circ (1 + \sigma)$.
Since $\sigma^2 = (\alpha^2 \otimes \alpha^2) \circ c^2$, the map $\sigma$ is an involution of $A \otimes A$ when $\alpha^2 = \id_A$ and $c^2 = \id_{A \otimes A}$, in particular in a symmetric category; it is a braiding only when $\alpha$ is compatible with the braided structure, which we do not assume.
For vector spaces with the swap and an involutive automorphism or anti-automorphism $\alpha$, the map $\sigma$ is the twisted swap of the first paper.
There the derivations of $\mu_\sigma$ that commute with $\alpha$ equal the derivations of $\mu$ that commute with $\alpha$, for a commutative reduced algebra with $2$ invertible \citep[Theorem 4.1]{Adrabi2026a}.
We prove the part of that theorem that holds in a general category.
We do not assume that $\alpha$ is compatible with the product.
Fuchs, Runkel and Schweigert use an endomorphism $\varsigma \colon A \to A$ with $\varsigma \circ \mu = \mu \circ c \circ (\varsigma \otimes \varsigma)$ \citep[Definition 2]{FuchsRunkelSchweigert2005}.
Their $\varsigma$ is a map on $A$, not on $A \otimes A$ like the $\sigma$ above.
We do not use that rule.

\begin{theorem}\label{thm:twisted}
Let $D \colon A \to A$ satisfy $D \alpha = \alpha D$.
\begin{enumerate}[label=(\roman*)]
\item $D \mu_\sigma - \mu_\sigma \Delta_D = E_D \circ (1 + \sigma)$.
\item If $\mu$ is commutative and $D$ is a derivation of $\mu_\sigma$, then $E_D + E_D \circ (\alpha \otimes \alpha) = 0$.
\item Let $\varepsilon \in \{1, -1\}$ and let $p, p' \colon X \to A$ satisfy $\alpha p = \varepsilon p$ and $\alpha p' = \varepsilon p'$. If $\mathcal{C}$ is $2$-torsion free, $\mu$ is commutative and $D$ is a derivation of $\mu_\sigma$, then $E_D \circ (p \otimes p') = 0$.
\end{enumerate}
\end{theorem}

\begin{proof}
(i) Since $D \alpha = \alpha D$, we have $\Delta_D (\alpha \otimes \alpha) = (\alpha \otimes \alpha) \Delta_D$.
With Lemma~\ref{lem:natural}, $\Delta_D \sigma = \sigma \Delta_D$.
The computation of Proposition~\ref{prop:defect-sym} gives (i).

(ii) By (i) and the hypothesis, $E_D (1 + \sigma) = 0$.
Now $E_D \sigma = E_D (\alpha \otimes \alpha) c$, and $c$ commutes with $\alpha \otimes \alpha$ because $c$ is natural.
As in the proof of Theorem~\ref{thm:commutative}, $E_D c = E_D$, so $E_D \sigma = E_D (\alpha \otimes \alpha)$.
Hence $E_D + E_D (\alpha \otimes \alpha) = 0$.

(iii) Since $(\alpha \otimes \alpha)(p \otimes p') = \varepsilon^2 (p \otimes p') = p \otimes p'$, composing (ii) with $p \otimes p'$ gives $E_D (p \otimes p') + E_D (p \otimes p') = 0$.
So $E_D (p \otimes p') = 0$.
\end{proof}

For vector spaces with $2$ invertible and $\alpha^2 = \id_A$, put $A_\pm = \ker(\alpha \mp 1)$; then $A = A_+ \oplus A_-$.
Part (iii) says that $E_D$ vanishes on $A_+ \otimes A_+$ and on $A_- \otimes A_-$. % claim-ok: restates Theorem 5.1(iii) for vector spaces
The vanishing on $A_+ \otimes A_-$ does not follow from (iii): it needs a reduced commutative algebra \citep[Theorem 4.1]{Adrabi2026a}. % claim-ok: a limit of Theorem 5.1; what [1, Theorem 4.1] adds is the hypotheses that close it

We end with a question.
In the survey of Bre\v{s}ar and Zelmanov, algebras are associative and over a field of characteristic not $2$ \citep[Section 1 and Section 11]{BresarZelmanov2025}, and the following theorem holds.
An algebra is \emph{semiprime} if $aAa = 0$ implies $a = 0$.

\begin{theorem}[{\citep[Theorem 11.1]{BresarZelmanov2025}}]\label{thm:semiprime}
A Jordan derivation of a semiprime associative algebra over a field of characteristic not $2$ is a derivation.
\end{theorem}

The theorem is due to Cusack \citep[Corollary 5]{Cusack1975}, as the survey notes.

For vector spaces with the flip as braiding, Corollary~\ref{cor:inclusion} gives $\Der(\mu) \subseteq \Der(\mu_+)$, and Theorem~\ref{thm:semiprime} gives the converse inclusion for semiprime associative algebras.
Which hypotheses on $(\mathcal{C}, A)$ give $\Der(\mu_+) = \Der(\mu)$ when $\mu$ is not commutative?

\begin{thebibliography}{99}

\bibitem[Adrabi(2026a)]{Adrabi2026a}
Adrabi, A.: Braided-symmetric derivations of algebras. Preprint (2026). \url{https://github.com/abderrahim-lectures/braided-sym-deriv}.

\bibitem[Adrabi(2026b)]{Adrabi2026b}
Adrabi, A.: Jordan and Lie derivations as derivations of the symmetrized and antisymmetrized products. Preprint (2026). \url{https://github.com/abderrahim-lectures/jordan-lie-symmetrized}.

\bibitem[Barnes et~al.(2016)]{BarnesSchenkelSzabo2016}
G.~E. Barnes, A. Schenkel and R.~J. Szabo, Nonassociative geometry in quasi-Hopf representation categories II: connections and curvature, \emph{J. Geom. Phys.} 106 (2016), 234--255.

\bibitem[Bre\v{s}ar(1988)]{Bresar1988}
Bre\v{s}ar, M.: Jordan derivations on semiprime rings. Proc. Amer. Math. Soc. \textbf{104}(4), 1003--1006 (1988).

\bibitem[Bre\v{s}ar and Zelmanov(2025)]{BresarZelmanov2025}
Bre\v{s}ar, M., Zelmanov, E.: Jordan homomorphisms: a survey. Serdica Math. J. \textbf{51}, 203--230 (2025).

\bibitem[Cusack(1975)]{Cusack1975}
Cusack, J.M.: Jordan derivations on rings. Proc. Amer. Math. Soc. \textbf{53}(2), 321--324 (1975).

\bibitem[Fuchs et~al.(2005)]{FuchsRunkelSchweigert2005}
Fuchs, J., Runkel, I., Schweigert, C.: Ribbon categories and (unoriented) CFT: Frobenius algebras, automorphisms, reversions. arXiv:math/0511590 (2005).

\bibitem[Graziani et~al.(2015)]{GrazianiMakhloufMeniniPanaite2015}
Graziani, G., Makhlouf, A., Menini, C., Panaite, F.: BiHom-associative algebras, BiHom-Lie algebras and BiHom-bialgebras. SIGMA \textbf{11}, 086 (2015).

\bibitem[Herstein(1957)]{Herstein1957}
Herstein, I.N.: Jordan derivations of prime rings. Proc. Amer. Math. Soc. \textbf{8}, 1104--1110 (1957).

\bibitem[Ion(2010)]{Ion2010}
Ion, B.: Relative PBW type theorems for symmetrically braided Hopf algebras. Preprint (2010). \url{https://arxiv.org/abs/1004.5022}.

\bibitem[Jacobson(1937)]{Jacobson1937}
Jacobson, N.: Abstract derivation and Lie algebras. Trans. Amer. Math. Soc. \textbf{42}(2), 206--224 (1937).

\bibitem[Jordan et~al.(1934)]{JordanVonNeumannWigner1934}
Jordan, P., von Neumann, J., Wigner, E.: On an algebraic generalization of the quantum mechanical formalism. Ann. of Math. \textbf{35}(1), 29--64 (1934).

\bibitem[Joyal and Street(1993)]{JoyalStreet1993}
Joyal, A., Street, R.: Braided tensor categories. Adv. Math. \textbf{102}(1), 20--78 (1993).

\bibitem[Martindale(1964)]{Martindale1964}
Martindale, W.S., III: Lie derivations of primitive rings. Michigan Math. J. \textbf{11}, 183--187 (1964).

\bibitem[Pareigis(1998)]{Pareigis1998}
Pareigis, B.: On Lie algebras in the category of Yetter--Drinfeld modules. Appl. Categ. Structures \textbf{6}, 151--175 (1998); arXiv:q-alg/9612023.

\end{thebibliography}

\end{document}
